%!TEX program = xelatex
\documentclass[11pt]{article}

\usepackage[margin=1in]{geometry}
\usepackage{amsmath, amssymb, amsthm}
\usepackage{graphicx}
\usepackage{hyperref}
\usepackage{xcolor}
\usepackage{booktabs}
\usepackage{multirow}
\usepackage{float}
\usepackage[section]{placeins}

\newcommand{\safeincludegraphics}[2][]{%
  \IfFileExists{#2}{\includegraphics[#1]{#2}}{%
    \fbox{\parbox[c][0.22\textheight][c]{0.9\linewidth}{\centering Missing figure:\\ \texttt{\detokenize{#2}}}}%
  }%
}

\newtheorem{theorem}{Theorem}
\newtheorem{assumption}{Assumption}

\setlength{\parindent}{0pt}
\setlength{\parskip}{1\baselineskip}

\title{Shared Spikes and the Conditional Optimality of Quadratic Routing in Anisotropic LLM Activations}
\author{Theory Group}
\date{\today}

\begin{document}
\maketitle

% {\color{red}
% \input{comments}
% }
% \newpage

\section{Recent Updates}

\paragraph{Current States}
\begin{enumerate}
    \item Under our modeling and assumptions of activations, through Bayesian Theories frame analysis, quadratic router is the best router for semantics rouitng.
    \item The quadratic router cannot work well at the very beginning of training, if we want to test its capability, we may need to add load balance warm-up procedures.
    \item  As for the specialization of experts, there's no existing definitions now, we may need to define one ourselves.
\end{enumerate}

\section{Background and Theoretical Critique}

Trained Large Language Models (LLMs) exhibit universal anisotropy in their representations. Severe anisotropy degrades routing in Mixture-of-Experts (MoE) architectures, because routing scores become dominated by shared components rather than token-specific semantic structure. 

To formalize this, our recent theoretical modeling postulates that the activation matrix $X$ can be decomposed into a shared spike, a semantic tail, and isotropic noise:
\begin{equation}
    x = c + u_k + \epsilon, \quad c \in C, \ u_k \in U_k
\end{equation}
where $C$ is a low-rank shared subspace (the spike) and $U_k$ is a feature-specific low-rank subspace (the semantic tail). 

\textbf{Critique of the Theoretical Proof:} The derivation of the Bayes-optimal quadratic router within this framework is mathematically rigorous. However, calculating the class-conditional covariance on $C^\perp$ introduces an implicit assumption. The formulation $\Sigma_{k,\perp} = \sigma_u^2 P_{U_k} + \tau^2(I - P_C)$ strictly requires $(I - P_C)P_{U_k}(I - P_C) = P_{U_k}$, which implies orthogonality between the semantic subspace $U_k$ and the spike subspace $C$ (i.e., $U_k \perp C$). While omitted in the original proof, this assumption is geometrically plausible in extreme high-dimensional spaces ($D \gg r$) due to concentration of measure. If we interpret the intermediate spectral components as the semantic tail space, this orthogonality assumption holds naturally.

\textbf{Clarification on Interventions:} Interventions targeting activation anisotropy should not conflate distinct operations. Mean removal is a targeted operation that eliminates the rank-1 DC offset and benefits linear routing projections. Empirically, mean removal yields a small but observable improvement in routing balance while maintaining expressive capacity; however, the precise mechanism still needs rigorous analysis. In contrast, full whitening is a global geometric transformation that rescales all covariances. Naive whitening can amplify isotropic noise $\epsilon$ and damage meaningful, non-isotropic organization. Interventions should therefore remain task-specific: isolate the shared mean channel from the semantic residual, rather than flatten the entire spectrum. Whitening is also computationally expensive, so deploying it in training should be treated as a last-resort option.

\section{Hypothesis and Mathematical Modeling}

\subsection*{1. The Mathematically Optimal Dispatch}
Our previous hypothesis posits that if token-specific semantics strictly reside within low-rank tail subspaces $U_k$, linear projections are fundamentally mis-specified. Under the assumption of symmetric Gaussian tails, the Bayes-optimal decision boundary is not a linear hyperplane, but rather a quadratic likelihood-ratio test on tail energies. 

The optimal routing policy is mathematically formulated as maximizing the projected energy within the respective orthogonalized semantic subspace:
\begin{equation}
    k^*(x) = \arg\max_k \|P_{U_k} (I - P_C)x\|^2
\end{equation}
This dictates that the optimal router must be a \textit{Spike-Invariant Degree-2 Router}. It first nullifies the global bias via the projection $(I-P_C)$ and subsequently evaluates the quadratic form of the residuals.

\subsection*{2. Computational Quantification in Subspaces}
Explicitly computing $\|P_{U_k} x_\perp\|^2$ requires storing and projecting onto $K$ distinct $D \times r$ basis matrices, yielding a routing complexity of $\mathcal{O}(K \cdot D \cdot r)$, which is computationally prohibitive. 

By operating entirely within a randomly projected subspace, we achieve substantial computational efficiency. Utilizing a random matrix $R \in \mathbb{R}^{D \times r}$ with entries drawn from $\mathcal{N}(0, 1/r)$, we approximate the preservation of subspace energy. The operation reduces to the following steps:

\begin{enumerate}
    \item \textbf{Spike Removal:} $x_\perp = x - P_C x$ ($\mathcal{O}(D \cdot r_c)$ FLOPs).
    \item \textbf{Random Projection:} $z = R^\top x_\perp$ ($\mathcal{O}(D \cdot r)$ FLOPs).
    \item \textbf{Quadratic Features:} $\phi(x) = [z; z \odot z] \in \mathbb{R}^{2r}$ ($\mathcal{O}(r)$ FLOPs).
    \item \textbf{Linear Dispatch:} $l(x) = W \phi(x)$ ($\mathcal{O}(E \cdot r)$ FLOPs).
\end{enumerate}

By truncating operations to subspace rank $r \ll D$ (e.g., $D=4096 \to r=64$), the routing complexity drops from $\mathcal{O}(D^2)$ to $\mathcal{O}(D \cdot r)$, making integration into the forward pass practical. However, random projection can degrade the signal-to-noise ratio. If the projection dimension $r$ is too small, subspace structure collapses, entangling semantic signals with isotropic noise.

\section{Analysis and Experimental Verification}

We tested the model assumptions on Qwen3-4B hidden activations from the Web of Science dataset, using four domains (Medical, Psychology, CS, biochemistry). Exp1--Exp3 use 1000 samples per domain; Exp4 uses 700/300 train-test splits, five random seeds, and layers $\{1,18,36\}$.

Following the notation of the provided gating note, we use
\begin{equation}
    x_\perp := (I-P_C)x, \qquad
    x_{c\perp} := (I-P_C)(x-\bar x_{\mathrm{train}}), \qquad
    z_0 := P_{C,1}x,
\end{equation}
where $P_C$ is the empirical spike projector, $P_{C,1}$ is the projector onto its first principal direction, and $\bar x_{\mathrm{train}}$ is the training mean used in Exp4. In the overlap plots, $U_A,U_B$ denote the top-$s$ tail bases of two domains, and the overlap score is
\begin{equation}
    \frac{\|U_A^\top U_B\|_F^2}{s}.
\end{equation}

Overall, the experiments support a \emph{refined} version of the theory: H2 is strongly supported, H4 is supported only in the deeper tail, and H3 is only partially supported if ``low-rank tail'' is interpreted literally.

\subsection*{Exp 1: a strong shared spike is present}

Table \ref{tab:exp1} shows that domain means are almost collinear across layers. At layer 12, both mean cosine and top-1 principal-direction cosine are essentially 1; at layers 4 and 36, the centered top-1 alignment drops, indicating that centering removes much of the universal direction.

\begin{table}[H]
\centering
\small
\caption{Exp1: layer-wise average alignment across the 6 domain pairs.}
\label{tab:exp1}
\begin{tabular}{@{}lccc@{}}
\toprule
\textbf{Layer} & \textbf{Mean Cosine} & \textbf{Centered PC1 Cosine} & \textbf{Uncentered PC1 Cosine} \\ \midrule
4  & 0.9884 & 0.6946 & 0.9884 \\
12 & 0.9960 & 0.9984 & 0.9964 \\
24 & 0.9921 & 0.9930 & 0.9924 \\
36 & 0.9509 & 0.7517 & 0.9506 \\ \bottomrule
\end{tabular}
\end{table}

\begin{figure}[H]
\centering
\begin{minipage}{0.43\linewidth}
\centering
\safeincludegraphics[width=\linewidth]{figures/exp1/pairwise_overlap_overlay_uncentered.png}
\end{minipage}\hfill
\begin{minipage}{0.43\linewidth}
\centering
\safeincludegraphics[width=\linewidth]{figures/exp1/pairwise_overlap_overlay_centered.png}
\end{minipage}
\caption{Exp1 at layer 12. Left: overlap before centering. Right: overlap after centering. The systematic reduction after centering indicates that a large fraction of the shared alignment comes from a common spike rather than class-specific tail structure.}
\label{fig:exp1}
\end{figure}

\subsection*{Exp 2: after spike removal, the residual spectrum is flatter but not truly low-rank}

After removing a rank-4 spike, the residual spectra of the four domains become very similar and much less spike-dominated (Figure \ref{fig:exp2}). However, the 90\% effective rank remains in the hundreds, especially in earlier layers. The data therefore do \emph{not} support a sharp finite-rank truncation of the residual; instead, they support a weaker statement: spike removal spreads the remaining energy across a broad residual spectrum.

\begin{table}[H]
\centering
\small
\caption{Exp2: 90\% effective rank of $x_\perp$ with spike rank $r_c=4$ (mean/min/max over the 4 domains).}
\label{tab:exp2}
\begin{tabular}{@{}lccc@{}}
\toprule
\textbf{Layer} & \textbf{Mean} & \textbf{Min} & \textbf{Max} \\ \midrule
1  & 374.50 & 373 & 376 \\
4  & 386.00 & 363 & 397 \\
12 & 330.00 & 322 & 338 \\
24 & 257.00 & 247 & 263 \\
36 & 263.75 & 245 & 275 \\ \bottomrule
\end{tabular}
\end{table}

\begin{figure}[H]
\centering
\safeincludegraphics[width=0.68\linewidth]{figures/exp2/scree_multidomain_rc4.png}
\caption{Exp2 at layer 24 with spike rank $r_c=4$. The residual spectra are nearly parallel across domains, which is consistent with a shared background geometry after removing the dominant spike; however, the slow decay shows that the residual is not strictly low-rank under a 90\% energy criterion.}
\label{fig:exp2}
\end{figure}

\subsection*{Exp 3: weak overlap appears mainly in the deeper tail}

The weak-overlap hypothesis does not hold uniformly across the full residual spectrum. It becomes more accurate when we ignore leading residual directions and focus on the deeper tail. At layer 12, the average overlap is only $0.0472$ for $(\texttt{tail\_start}=64, r=16)$, but rises to $0.1390$ when $r=64$. Pairwise structure is also non-uniform: \emph{Medical vs Psychology} is consistently more overlapping than \emph{Psychology vs CS}. For example, at layer 12 and \texttt{tail\_start}=64, the overlap is $0.0679$ vs. $0.0429$ at $r=16$, and $0.1823$ vs. $0.1247$ at $r=64$. Layer 24 is less orthogonal overall, indicating that H4 is layer-dependent rather than universal.

\begin{table}[H]
\centering
\small
\caption{Exp3: aggregated overlap $\|U_A^\top U_B\|_F^2/r$ for the deeper tail (\texttt{tail\_start}=64).}
\label{tab:exp3}
\begin{tabular}{@{}lcc@{}}
\toprule
\textbf{Layer} & \textbf{$r=16$} & \textbf{$r=64$} \\ \midrule
4  & 0.0625 & 0.1748 \\
12 & 0.0472 & 0.1390 \\
24 & 0.0879 & 0.2298 \\
36 & 0.0500 & 0.1363 \\ \bottomrule
\end{tabular}
\end{table}

\begin{figure}[H]
\centering
\safeincludegraphics[width=0.74\linewidth]{figures/exp3/pair_overlap_focus_layer12.png}
\caption{Selected pair comparison in the subspace-overlap analysis at layer 12. In both the mid-tail (\texttt{start}=32) and deeper tail (\texttt{start}=64), \emph{Medical vs Psychology} remains more overlapping than \emph{Psychology vs CS}, showing that the residual geometry is not uniformly incoherent across semantic pairs.}
\label{fig:exp3}
\end{figure}

\subsection*{Exp 4: quadratic routing is conditionally validated on the spike-removed residual}

The clearest evidence for the theorem appears only after separating the spike and residual channels. On raw pooled activations, a linear probe already reaches $0.9075 \pm 0.0584$ mean accuracy, so the full representation still contains a strong linear shortcut. In contrast, once we restrict to the centered residual $X_{c\perp}$, linear probing collapses to $0.5868 \pm 0.0508$, while the quadratic rule and the direct residual-energy rule recover $0.8940 \pm 0.0609$ and $0.8944 \pm 0.0609$, respectively. The experiments therefore do \emph{not} show that quadratic routing beats every baseline on raw activations; they show something more precise: \emph{inside the spike-removed residual space, the correct decision form is much closer to quadratic than to linear.}

\begin{table}[H]
\centering
\small
\caption{Exp4: selected routing baselines, averaged over all layers, seeds, and domain pairs.}
\label{tab:exp4}
\begin{tabular}{@{}lcc@{}}
\toprule
\textbf{Method} & \textbf{Mean Acc.} & \textbf{Std.} \\ \midrule
Linear LR on $X$ & 0.9075 & 0.0584 \\
Linear LR on $z_0$ & 0.7191 & 0.1441 \\
Linear LR on $X_{c\perp}$ & 0.5868 & 0.0508 \\
Quadratic-LR on $X_{c\perp}$ & 0.8940 & 0.0609 \\
Residual-Energy Rule on $X_{c\perp}$ & 0.8944 & 0.0609 \\ \bottomrule
\end{tabular}
\end{table}

\begin{figure}[H]
\centering
\safeincludegraphics[width=0.72\linewidth]{figures/exp4/router_selected_methods.png}
\caption{Exp4: selected baselines. $X$ is the raw pooled activation, $z_0$ is the first spike coordinate, and $X_{c\perp}=(I-P_C)(X-\bar x_{\mathrm{train}})$. The essential contrast is that linear routing fails on $X_{c\perp}$, whereas quadratic tail-energy routing almost fully restores separability.}
\label{fig:exp4}
\end{figure}

\section{Conclusion}

\noindent
\fcolorbox{black}{gray!10}{%
\parbox{0.97\linewidth}{%
\textbf{Takeaway} The experiments above provide strong evidence for a universal shared spike, partial evidence for weakly overlapping semantic tails, and a clear \emph{conditional} validation of the quadratic router: once the spike is removed, linear routing on the residual fails, while quadratic tail-energy routing recovers almost all of the separability. Under the current generative model, the spike-invariant quadratic rule is therefore the right static classifier, but real LLM routing must ultimately be judged by both \emph{representation geometry} and \emph{optimization dynamics}.}}

The strongest claim supported by the data is therefore not ``the entire tail is low-rank and orthogonal,'' but rather: there exists a dominant common spike, deeper residual directions are substantially more class-specific, and a degree-2 energy rule is the correct inductive bias on that spike-removed residual space. This yields: the theory is validated as a \emph{static} statement on frozen activations and the assumption may be too strong, while the training-time failure identifies where the present model is incomplete.

\section{Update and Next Iteration}

We also implemented a quadratic router and evaluated it as an actual routing module. In the current training setup, it exhibited substantial load imbalance: a small number of experts tended to dominate early, and the routing pattern became increasingly polarized over time. At this stage, we do not interpret this as evidence against the quadratic principle itself. Under the static generative model, the quadratic rule remains the natural Bayes-optimal decision rule. What the current result shows is more limited but also more important: \emph{static optimality alone is not sufficient to guarantee trainability or stable load distribution.}

This leads to a necessary update of the modeling perspective. The previous decomposition
\[
x = c + u_k + \epsilon
\]
is useful for describing representation geometry and for deriving an optimal decision rule on frozen features, but it is not yet sufficient for explaining routing behavior during training. The next iteration should therefore move to a geometry-plus-dynamics model, in which routing is characterized not only by the structure of \(x\), but also by the evolution of router scores, expert occupancy, and score dispersion over training.

A more cautious updated hypothesis is the following: the quadratic residual rule may still be the correct \emph{static} target under the spike-tail model, but its practical success depends on whether its score scale, competition pattern, and expert loads remain sufficiently controlled during optimization. The next step is therefore not to abandon the quadratic direction, but to study the conditions under which it can be trained stably. Concretely, this motivates analyzing the gradient scale and score dispersion of the quadratic gate, and then testing simple stabilization mechanisms such as scale-controlled initialization, clipping, or explicit balancing terms.

\end{document}
